% mtheory_duality.tex -- Physical Review D submission
\documentclass[aps,prd,twocolumn,showpacs,preprintnumbers,superscriptaddress,floatfix]{revtex4-2}

\usepackage{amsmath,amssymb,amsfonts}
\usepackage{graphicx}
\usepackage{hyperref}
\usepackage{booktabs}

\newcommand{\SU}{\mathrm{SU}}
\newcommand{\UU}{\mathrm{U}}
\newcommand{\SO}{\mathrm{SO}}
\newcommand{\Sp}{\mathrm{Sp}}
\newcommand{\Nf}{N_{\!f}}
\newcommand{\Nc}{N_{\!c}}
\newcommand{\Tr}{\mathrm{Tr}}
\newcommand{\adj}{\mathrm{adj}}
\newcommand{\diag}{\mathrm{diag}}
\newcommand{\keV}{\,\mathrm{keV}}
\newcommand{\MeV}{\,\mathrm{MeV}}
\newcommand{\GeV}{\,\mathrm{GeV}}
\newcommand{\TeV}{\,\mathrm{TeV}}
\newcommand{\ie}{i.e.\@}
\newcommand{\eg}{e.g.\@}
\newcommand{\zch}{\mathfrak{z}}
\newcommand{\dd}{\mathrm{d}}

\begin{document}

\title{Dirac--Majorana duality from the M-theory lift
  of the sBootstrap}

\author{A.~Rivero}
\email{rivero@unizar.es}
\affiliation{Departamento de F\'{\i}sica Te\'orica,
  Universidad de Zaragoza, 50009 Zaragoza, Spain}

\date{\today}

\begin{abstract}
The sBootstrap identifies Standard Model leptons with mesino
composites of a confining $\SU(3)_F$ family gauge theory at
$\Nf = \Nc = 3$.
In the Hanany--Witten brane construction, mesons lift to
M2-branes coupling to the 3-form~$C_3$, while baryons lift to
M5-branes coupling to the 6-form~$C_6$.
We show that this M2/M5 correspondence maps onto the distinction
between Dirac and Majorana mass types: the bilinear meson
$M = Q\tilde{Q}$ produces Hermitian (Dirac) mass matrices, while
the trilinear baryon $B = \varepsilon Q^3$ produces symmetric
(Majorana) mass matrices, with the involution
$M_i \leftrightarrow \adj(M)_i$ lifting to $C_3 \leftrightarrow C_6$
electromagnetic duality.
The meson K\"ahler metric is approximately diagonal (protecting
small CKM angles), whereas the baryon K\"ahler metric is
generically democratic (allowing large PMNS angles).
The holomorphic superpotential $W = BM\tilde{B}/\Lambda^5$
generates a type-I seesaw structure with
$M_R^{ii} = m_i$ (exact from $s$-confinement); however, the
physical neutrino mass scale requires non-perturbatively large
K\"ahler suppression factors of order~$10^5$, which cannot be
computed from holomorphy alone.
Using multi-loop SM running of Yukawa couplings, the inverse
Koide ratio $K^{-1}(d,s,b)$ crosses~$2/3$ at
$Q \approx 10\TeV$ in the Standard Model with no new thresholds,
while in the MSSM it remains $0.23$--$0.65\%$
above~$2/3$ at~$M_{\rm GUT}$ for all~$\tan\beta$---disfavouring
conventional low-energy supersymmetry.
\end{abstract}

\maketitle


%======================================================================
\section{Introduction}
\label{sec:intro}
%======================================================================

The sBootstrap~\cite{Rivero:2005b,Rivero:2024epjc} is a structural
effective field theory that identifies Standard Model fermions with
composites of a confining $\SU(3)_F$ family gauge theory at
$\Nf = \Nc = 3$.
The nine mesinos of the meson matrix
$M^i{}_j = Q^i\tilde{Q}_j$ decompose under the electroweak
branching $\mathbf{3}_L \to \mathbf{2}_{-1/2} \oplus
\mathbf{1}_{+1}$ into three lepton doublets~$L_i$ and three
charged-lepton singlets~$e^c_i$.
Quarks, which carry colour and therefore cannot be family-singlet
composites, enter through partial
compositeness~\cite{Kaplan:1991dc}.
The adjugate $\adj(M) = \det(M)\,M^{-1}$ provides the second
MSSM Higgs doublet, and the charge-squared mass law
$m_i = \mu\,\zch_i^2$ yields the Koide relation
$K = 2/3$ as a trace identity of the canonical
$\UU(3)_F$ normalisation~\cite{Rivero:2024epjc}.

The Seiberg seesaw~\cite{Seiberg:1994bz} produces a direct branch
($\langle M_i \rangle \propto 1/m_i$) and an inverse branch
($\adj(M)_i \propto m_i$).
In the M-theory lift of the Hanany--Witten
construction~\cite{HananyWitten:1997}, mesons are M2-branes
stretched between separated parts of the M5 worldvolume curve,
coupling to the 3-form potential~$C_3$; baryons are M5-branes
wrapping compact 4-cycles, coupling to the
6-form~$C_6$~\cite{Witten:1997M5,Witten:1998baryon}.
The Hodge duality $*G_4 = G_7$ in 11-dimensional
supergravity~\cite{CremmerJuliaScherk:1978} is the field-strength
avatar of M2/M5 electromagnetic duality.

In this paper we develop the observation that this M2/M5
correspondence maps onto the distinction between Dirac and Majorana
mass types.
Specifically:
(i)~the bilinear meson $M = Q\tilde{Q}$ produces Hermitian (Dirac)
mass matrices, while the trilinear baryon $B = \varepsilon Q^3$
produces symmetric (Majorana) mass matrices;
(ii)~the meson K\"ahler metric is approximately diagonal (small CKM
angles), while the baryon K\"ahler metric is generically democratic
(large PMNS angles);
(iii)~the holomorphic superpotential
$W = BM\tilde{B}/\Lambda^5$ generates a type-I seesaw structure.

We also present updated renormalisation group evidence for the
inverse Koide relation in the down-quark sector, using the
multi-loop running Yukawa couplings of Antusch, Hinze, and
Saad (AHS)~\cite{AHS:2024}, and show that the SM desert
uniquely produces the crossing $K^{-1}(d,s,b) = 2/3$ at
$Q \approx 10\TeV$, while the MSSM fails to reach this value.

The paper is organised as follows.
Section~\ref{sec:dictionary} establishes the brane dictionary at
each stage of the breaking chain.
Section~\ref{sec:duality} develops the $C_3 \leftrightarrow C_6$
Dirac--Majorana duality.
Section~\ref{sec:kahler} discusses the K\"ahler protection
asymmetry.
Section~\ref{sec:seesaw} derives the geometric seesaw and
honestly assesses its limitations.
Section~\ref{sec:RG} presents the RG evidence for the SM desert.
Section~\ref{sec:discussion} discusses implications, including
the debunking of a previously noted coincidence with the
Pendleton--Ross fixed point.
Section~\ref{sec:conclusions} summarises and identifies
experimental tests.

Throughout, we work in natural units ($\hbar = c = 1$)
with metric signature $(-,+,+,+)$ and generators normalised as
$\Tr(T^a T^b) = \frac{1}{2}\delta^{ab}$.
The 11d supergravity conventions follow
Cremmer--Julia--Scherk~\cite{CremmerJuliaScherk:1978}: $C_3$ is
the 3-form potential, $G_4 = \dd C_3$, $G_7 = *G_4 = \dd C_6$.


%======================================================================
\section{The brane dictionary}
\label{sec:dictionary}
%======================================================================

We first establish the M-theory lift of each sBootstrap composite.
The type~IIA Hanany--Witten construction for $\SU(3)$ SQCD with
$\Nf = 3$ places $\Nc = 3$ D4-branes suspended between two
NS5-branes, with three D6-branes providing fundamental
flavours~\cite{HananyWitten:1997}.
Lifting to M-theory by opening the $x^{10}$
circle~\cite{Witten:1997M5}:
the D4-branes become an M5-brane wrapping a holomorphic curve
$\Sigma$ in the $(x^4, x^5, x^6, x^{10})$ space;
the D6-branes become Taub-NUT centres; the NS5-branes remain
as~M5.

The composites have direct brane interpretations, collected in
Table~\ref{tab:dictionary}.

\begin{table}
\centering\small
\caption{M-theory dictionary for $\SU(3)$ SQCD with $\Nf = \Nc = 3$.
  The column ``$C$-field'' indicates the dominant coupling.
  The mass-type column identifies the fermion bilinear structure
  in the low-energy effective theory.}
\label{tab:dictionary}
\begin{tabular}{llllll}
\toprule
Composite & IIA & M-theory & $C$-field & Mass type & $p$ \\
\midrule
$M = Q\tilde{Q}$ & D6--D6 & M2 & $C_3$ & Dirac & 2 \\
$B = \varepsilon Q^3$ & wrapped D4 & wrapped M5 & $C_6$ & Majorana & 5 \\
$\adj(M)$ & thr.\ D4 & thr.\ M5 & mixed & inv.\ branch & --- \\
$M \leftrightarrow \adj(M)$ & T-dual & EM dual & $C_3{\leftrightarrow}C_6$ & --- & --- \\
$\Phi$ (aux.\ adj.) & D4--D4 & M5 modulus & metric & --- & --- \\
\bottomrule
\end{tabular}
\end{table}

\paragraph{M2-branes and mesons.}
The meson $M^i{}_j = Q^i\tilde{Q}_j$ lifts to an M2-brane
ending on the M5 worldvolume.
The M2 worldvolume action contains the Wess--Zumino term
$\int C_3$, so the meson couples electrically to~$C_3$.
The M2-brane area encodes the meson mass: for a $p$-brane
whose worldvolume dimensions scale with the family
charge~$\zch_i$, the mass scales as
$m_i \propto \zch_i^p$~\cite{Rivero:2024epjc}.
The generalised Koide ratio
$K_p = \sum \zch_i^{2p}/(\sum \zch_i^p)^2$ is independent of the
Cartan direction only for $p = 2$---the M2-brane.
This consistency check does not constitute a derivation of the
charge-squared law, since assuming both worldvolume directions
scale with $\zch_i$ is equivalent to the bilinear structure
$M = Q\tilde{Q}$.

\paragraph{M5-branes and baryons.}
The baryon $B = \varepsilon_{ijk}Q^iQ^jQ^k$ lifts to an M5-brane
wrapping $S^4$ in the near-horizon geometry of the D6-branes,
with three string
endpoints~\cite{Witten:1998baryon,HoriOoguriOz:1997}.
The M5 couples electrically to $C_6$.
For $p = 5$, the ratio $K_5$ depends on the Cartan angle;
baryon masses are determined algebraically
($M_R^{ii} = m_i$, exact from $s$-confinement~\cite{Seiberg:1994bz})
rather than by a power law in~$\zch_i$.
The baryon mass is a topological quantity (the volume of $S^4$
in the Taub-NUT geometry), while the meson mass is a modular
quantity (the area of a 2-cycle on~$\Sigma$).

\paragraph{The adjugate.}
The adjugate $\adj(M)_i = \varepsilon_{ijk} M_j M_k / \Lambda^3$
involves the $\varepsilon$ tensor but is a product of two mesons
with an $\varepsilon$ contraction.
In M-theory, this corresponds to three M2-branes joined at a
baryon-like vertex---a mixed $C_3/C_6$ configuration, not a pure
$C_6$ coupling.
Only $B = \varepsilon Q^3$ is a pure $C_6$ object.
The involution $M_i \leftrightarrow \adj(M)_i$ exchanges pure
$C_3$ with mixed $C_3/C_6$, not pure $C_3$ with pure~$C_6$.

\paragraph{The breaking chain.}
The full correspondence across the breaking chain
$\text{M-theory} \to \mathcal{N}{=}2 \to \mathcal{N}{=}1 \to
\text{SM}$ is:
\begin{itemize}
\item The M5-brane wrapping $\Sigma$ gives $\mathcal{N}{=}2$ in
  4d; the genus-2 curve encodes the Seiberg--Witten
  prepotential~\cite{SeibergWitten:1994}.
\item The adjoint mass $\mu_\Phi$ deforms $\Sigma$, breaking
  $\mathcal{N}{=}2 \to \mathcal{N}{=}1$.
  The hypermultiplet $(M_i, \adj(M)_i)$ splits into two
  independent $\mathcal{N}{=}1$ chirals---the two periods
  of~$\Sigma$.
\item At confinement, the 9~mesinos identify with the SM
  lepton content: $3L_i + 3e^c_i$.
  Baryoninos provide sterile (Majorana) neutrinos.
\end{itemize}

\paragraph{Colour-3 impossibility.}
The 11d 3-form decomposes as
$C_3 \in \Lambda^3(\mathbb{R}^9) = \mathbf{84}$, which under
$\SU(3)_c$ contains $\mathbf{1}_c$, $\mathbf{3}_c$,
$\overline{\mathbf{3}}_c$ but never $\mathbf{6}_c$.
The colour-sextet representation appears only in the symmetric
(gravitational) sector
$g_{MN} \in \mathrm{Sym}^2(\mathbb{R}^9) = \mathbf{45}$.
This is consistent with the field-theory result that
$\SU(3)_F$ composites with $\Nf = \Nc = 3$ never carry the
colour representation required for quarks to be fully
composite~\cite{Rivero:2024epjc}.


%======================================================================
\section{Dirac--Majorana duality}
\label{sec:duality}
%======================================================================

We now develop the central observation: the $C_3/C_6$ distinction
maps onto the Dirac/Majorana mass distinction.

\subsection{Electric and magnetic mass types}

In 11d supergravity, the 4-form field strength
$G_4 = \dd C_3$ and its Hodge dual $G_7 = *G_4 = \dd C_6$
satisfy~\cite{CremmerJuliaScherk:1978,Townsend:1995}:
\begin{align}
  \dd * G_4 &= *j_3 \quad (\text{M2 source, electric})\,,
  \label{eq:M2source} \\
  \dd G_4 &= j_6 \quad
  (\text{M5 source, magnetic})\,.
  \label{eq:M5source}
\end{align}
The M2-brane is the electric source (appearing on the right-hand
side of $\dd*F = j$); the M5-brane is the magnetic source
(modifying the Bianchi identity).

In the sBootstrap, the algebraic structure of each composite
determines the symmetry of its mass matrix:
\begin{itemize}
\item \textbf{Meson} $M^i{}_j = Q^i\tilde{Q}_j$ (bilinear):
  the mass matrix is Hermitian, coupling left-handed to
  right-handed fermions---Dirac type.
\item \textbf{Baryon} $B = \varepsilon_{ijk}Q^iQ^jQ^k$ (trilinear):
  the mass matrix is symmetric, coupling a fermion to itself and
  violating lepton number by two units---Majorana type.
\end{itemize}
This is an algebraic fact of
$s$-confinement~\cite{Seiberg:1994bz,CsakiSchmaltzSkiba:1997},
independent of M-theory.
The M-theory lift upgrades this field-theory observation to a
geometric statement: the Dirac/Majorana distinction is the
electric/magnetic distinction of the
$C$-field~(Table~\ref{tab:duality}).

\begin{table}
\centering\small
\caption{Structural correspondence between $C$-field sectors and
  fermion mass types.  The identification is algebraic at the
  field-theory level and geometric in M-theory.}
\label{tab:duality}
\begin{tabular}{lll}
\toprule
Property & Electric ($C_3$, M2) & Magnetic ($C_6$, M5) \\
\midrule
Brane & M2 & M5 \\
Composite & Meson $Q\tilde{Q}$ & Baryon $\varepsilon Q^3$ \\
Mass type & Dirac (bilinear) & Majorana (trilinear) \\
Matrix symmetry & Hermitian & Symmetric \\
Protected by & Chiral symmetry & Lepton number \\
SM sector & Charged fermions & Neutrinos \\
\bottomrule
\end{tabular}
\end{table}

\subsection{The involution and A/B period duality}

The Seiberg--Witten curve $\Sigma$ for $\SU(3)$ with $\Nf = 3$ is
a genus-2 hyperelliptic
surface~\cite{SeibergWitten:1994,ArgyresPlesserShapere:1996}:
\begin{equation}
  y^2 = \bigl[x^3 - u_2\,x - u_3\bigr]^2
    - 4\Lambda^3\prod_{i=1}^{3}(x + m_i)\,,
  \label{eq:SWcurve}
\end{equation}
where $u_2 = \Tr(\Phi^2)$ and $u_3 = \det(\Phi)$ are Coulomb-branch
moduli.
In the confining vacuum, these are fixed by the fermion mass
matrix $D = \diag(m_1, m_2, m_3)$ via the Cayley--Hamilton
equation of the $\mathcal{N}{=}2$ adjoint~$\Phi$:
\begin{equation}
  u_2 = \Tr(D^2) = \sum_i m_i^2\,,\quad
  u_3 = 2\det(D) = 2\prod_i m_i\,.
  \label{eq:Coulomb_moduli}
\end{equation}
The four periods of $\Sigma$---two $A$-periods
$a_i = \oint_{A_i}\lambda_{\rm SW}$ and two $B$-periods
$a_{D,i} = \oint_{B_i}\lambda_{\rm SW}$---encode both mass
hierarchies:
\begin{align}
  A\text{-periods:}\quad &\langle M_i \rangle
    \propto 1/m_i
    \quad (\text{direct branch})\,,
    \label{eq:direct} \\
  B\text{-periods:}\quad &\adj(M)_i
    \propto m_i
    \quad (\text{inverse branch})\,.
    \label{eq:inverse}
\end{align}
The involution $M_i \leftrightarrow \adj(M)_i$ exchanges
$A$-periods with $B$-periods.
In M-theory, this is precisely the electromagnetic duality
$C_3 \leftrightarrow C_6$, exchanging M2 and M5
brane data.

\subsection{Scope and caveats}

Two important qualifications:

(i)~The $C_3 \leftrightarrow C_6$ duality is a duality
\emph{within} the mesonic sector (direct $\leftrightarrow$ inverse
branch), not a duality \emph{between} the mesonic and baryonic
sectors.
The baryon $B = \varepsilon Q^3$ wraps $S^4$, not~$\Sigma$; it is
a topologically distinct object from $\adj(M)$.
The connection between mesons and baryons is mediated by the
quantum-modified constraint
\begin{equation}
  \det M - B\tilde{B} = \Lambda^6\,,
  \label{eq:quantum_constraint}
\end{equation}
not by a direct $C_3 \leftrightarrow C_6$ map.

(ii)~The correspondence in Table~\ref{tab:duality} is structural,
not quantitative.
No formula relating the Dirac mass of a given fermion to its
Majorana partner through the duality currently exists.
The most promising path to such a relation is through the Riemann
bilinear relations of the genus-2 period matrix, which constrain
the product of $A$- and $B$-periods; this computation has not
been carried out for the specific $\SU(3)$, $\Nf = 3$ case at
the Koide moduli point.


%======================================================================
\section{K\"ahler protection asymmetry}
\label{sec:kahler}
%======================================================================

In the $s$-confined theory, the holomorphic masses of both mesinos
and baryoninos are protected by the
$\mathcal{N}{=}2$ non-renormalisation theorem (the Higgs-branch
hyper-K\"ahler metric is one-loop exact in the $\mathcal{N}{=}2$
limit~\cite{SeibergWitten:1994}).
However, the physical masses
$m_{\rm phys} = m_{\rm hol}/\sqrt{Z}$ depend on the K\"ahler
metric~$Z$, which is not protected by holomorphy after
$\mathcal{N}{=}2 \to \mathcal{N}{=}1$ breaking.
The key structural difference between the two sectors is in the
flavour structure of~$Z$.

\paragraph{Meson K\"ahler metric: approximately diagonal.}
The meson $M^i{}_j$ carries definite flavour quantum numbers
$(i, j)$.
In the $\mathcal{N}{=}2$ limit, the K\"ahler metric on the Higgs
branch is exactly diagonal in the flavour
basis~\cite{SeibergWitten:1994}.
Breaking to $\mathcal{N}{=}1$ introduces corrections controlled by
\begin{equation}
  \epsilon = \frac{m_e}{m_\tau} \approx 2.9 \times 10^{-4}\,,
  \label{eq:epsilon}
\end{equation}
giving $\delta Z_M / Z_M \sim \epsilon^2 \sim 10^{-7}$
to~$10^{-9}$~\cite{Rivero:2024epjc}.
This is consistent with the observed precision of
$K(e,\mu,\tau) = 2/3$ to~$0.001\%$.
The near-diagonality of $Z_M$ means that the meson-sector mass
matrix remains approximately diagonal after K\"ahler rotation:
\emph{the CKM matrix is small because meson flavour numbers
protect the K\"ahler metric.}

\paragraph{Baryon K\"ahler metric: generically democratic.}
The baryon $B = \varepsilon_{ijk}Q^iQ^jQ^k$ involves all three
flavours symmetrically through the $\varepsilon$ tensor.
No flavour quantum number distinguishes $B$ from any
permutation of its constituent quarks.
Consequently, the baryon K\"ahler metric $Z_B^{ij}$ has no reason
to be diagonal: it is generically of democratic form, with
$\mathcal{O}(1)$ off-diagonal entries.
The K\"ahler rotation needed to bring the physical mass matrix to
canonical form then introduces large mixing angles.
\emph{The PMNS matrix is large because the $\varepsilon$ tensor
democratises the baryon K\"ahler metric.}

This is a structural explanation, not a numerical prediction.
Computing $Z_B$ requires the non-perturbative K\"ahler potential
at strong coupling, which is beyond current techniques.
Nevertheless, the qualitative asymmetry---diagonal mesonic vs.\
democratic baryonic---is a robust consequence of the composite
structure and does not depend on the specific form of the
K\"ahler potential.
This asymmetry has been noted
previously~\cite{MasieroVeneziano:1985,Karasik:2020}
in related compositeness frameworks;
the present work identifies the M-theory origin as the
$C_3$ vs.\ $C_6$ distinction.


%======================================================================
\section{The geometric seesaw}
\label{sec:seesaw}
%======================================================================

\subsection{Holomorphic superpotential}

The full holomorphic superpotential on the mesonic branch of
$\SU(3)_F$ with $\Nf = \Nc = 3$
is~\cite{Seiberg:1994bz,CsakiSchmaltzSkiba:1997}:
\begin{equation}
  W = X\bigl(\det M - B\tilde{B} - \Lambda^6\bigr)
    + \Tr(m\,M) + \frac{1}{\Lambda^5}\,B\,M\,\tilde{B}\,,
  \label{eq:superpotential}
\end{equation}
where $X$ is a Lagrange multiplier enforcing the
quantum-modified constraint, $m = \diag(m_1, m_2, m_3)$ are the
UV mass deformations, and the last term is the baryon--meson
coupling (dimension~8, suppressed by $\Lambda^5$).

\subsection{Baryonino Majorana masses}

On the mesonic branch ($B = \tilde{B} = 0$), the $F$-term
equations give the standard
result~\cite{Seiberg:1994bz,Rivero:2024epjc}:
\begin{equation}
  \langle M_i \rangle = \frac{\Lambda^2 P^{1/3}}{m_i}\,,
  \qquad P = m_1 m_2 m_3\,,
  \label{eq:seiberg_seesaw}
\end{equation}
where $P^{1/3} = (m_e m_\mu m_\tau)^{1/3} = 45.78\MeV$ for the
lepton sector.
The baryonino masses follow from the second derivative of~$W$
with respect to the baryonic fields.
The exact $s$-confinement result is~\cite{Rivero:2024epjc}:
\begin{equation}
  M_R^{ii} = m_i\,.
  \label{eq:baryonino_mass}
\end{equation}
That is, the Majorana mass of the $i$-th baryonino equals the UV
mass deformation.
For the lepton sector:
$M_R = \diag(m_e, m_\mu, m_\tau)
= \diag(0.511, 105.66, 1776.86)\MeV$.

\subsection{Dirac mass from the baryon--meson coupling}

The third term in Eq.~\eqref{eq:superpotential} generates a Dirac
mass for the baryonino when the meson~$M$ acquires a VEV.
After canonical normalisation of the confined fields
(K\"ahler metric $Z_M$ for mesons, $Z_B$ and $Z_{\tilde{B}}$ for
baryons), the physical Dirac mass is:
\begin{equation}
  m_D^i
    = \frac{\sqrt{Z_M}}{\sqrt{Z_B Z_{\tilde{B}}}}\,
      \frac{\Lambda\,P^{1/3}}{m_i}\,.
  \label{eq:dirac_mass}
\end{equation}
Setting all K\"ahler factors to unity gives the holomorphic
estimate $m_D^i = \Lambda P^{1/3}/m_i$.

\subsection{Type-I seesaw structure}

Integrating out the heavy baryonino at mass
$M_R^i = m_i$~[Eq.~\eqref{eq:baryonino_mass}], the light
neutrino mass matrix is:
\begin{equation}
  m_\nu^{ii}
    = \frac{(m_D^i)^2}{M_R^i}
    = \frac{Z_M}{Z_B Z_{\tilde{B}}}\,
      \frac{\Lambda^2 P^{2/3}}{m_i^3}\,.
  \label{eq:seesaw}
\end{equation}
This has the structure of a type-I seesaw,
$m_\nu = m_D^T M_R^{-1} m_D$, with the specific identifications:
$m_D^i$ from the meson VEV (electric/$C_3$ sector) and
$M_R^i = m_i$ from the baryonino F-term
(magnetic/$C_6$ sector).

On the M5 worldvolume, this coupling corresponds to the
interaction $\int_{M5} C_3 \wedge H_3$, where $H_3$ is the
self-dual 3-form on the M5 worldvolume.
The seesaw suppression $m_D^2/M_R$ arises because the M2
contributions (area-suppressed) combine with the M5 contribution
(volume-enhanced).
This identification is schematic; a rigorous derivation from the
M5 worldvolume effective action for $\SU(3)$, $\Nf = 3$ has not
been performed.

\subsection{The scale problem: an honest assessment}

The holomorphic seesaw~\eqref{eq:seesaw} at unit K\"ahler factors
gives $m_\nu^{ii} \propto 1/m_i^3$, predicting an extreme inverted
hierarchy ($m_{\nu_e} \gg m_{\nu_\tau}$) and masses in the GeV
range---wrong by $10$--$17$ orders of magnitude.
We collect the numerical values in Table~\ref{tab:seesaw}.

\begin{table}
\centering\small
\caption{Holomorphic seesaw masses at unit K\"ahler factors
  ($Z_M = Z_B = Z_{\tilde{B}} = 1$, $\Lambda = 100\GeV$),
  compared with observed neutrino mass estimates.
  The naive seesaw fails on both scale and hierarchy.}
\label{tab:seesaw}
\begin{tabular}{lccc}
\toprule
Generation & $m_\nu^{\rm hol}$ & $m_\nu^{\rm obs}$ (est.) & Ratio \\
\midrule
$e$   & $\sim 10^{11}\GeV$ & $\sim 0.35$ meV & $4.5 \times 10^{17}$ \\
$\mu$ & $\sim 10^{4}\GeV$  & $\sim 8.6$ meV  & $2.1 \times 10^{12}$ \\
$\tau$& $\sim 3.7\GeV$     & $\sim 50$ meV   & $7.5 \times 10^{10}$ \\
\bottomrule
\end{tabular}
\end{table}

To obtain $m_{\nu_3} \sim 50$~meV from the tau sector requires
$\sqrt{Z_B Z_{\tilde{B}}} \sim 2.7 \times 10^5$---an enormous
K\"ahler suppression that cannot be computed from holomorphy.
Furthermore, the tree-level PMNS matrix on the mesonic branch is
the identity: all neutrino masses are diagonal and there is no
lepton mixing.
Non-trivial PMNS angles require either off-diagonal K\"ahler
rotations in the baryonic sector, radiative corrections, or
SUSY-breaking baryon condensates~\cite{MasieroVeneziano:1985}.

\paragraph{Summary.}
The geometric seesaw provides the correct \emph{structure}:
Dirac masses from the meson VEV ($C_3$ sector), Majorana masses
from baryonino F-terms ($C_6$ sector), and a type-I seesaw formula.
It correctly identifies baryoninos as sterile neutrinos and the
meson VEV as the Dirac mass source.
It does \emph{not} predict the neutrino mass scale or PMNS mixing
without additional non-holomorphic input.
The framework gives the right \emph{players} but not the right
\emph{numbers}.


%======================================================================
\section{SM desert vs.\ MSSM: Koide RG evidence}
\label{sec:RG}
%======================================================================

The charge-squared law $m_i = \mu\,\zch_i^2$ and the canonical
$\UU(3)_F$ normalisation give the Koide ratio $K = 2/3$ as a trace
identity~\cite{Rivero:2024epjc}.
For the down-quark sector, which arises from the adjugate (inverse)
branch, the relevant observable is the inverse Koide
ratio~\cite{Rivero:2024epjc}:
\begin{equation}
  K^{-1}(d,s,b) \equiv
  \frac{1/m_d + 1/m_s + 1/m_b}
       {\bigl(1/\sqrt{m_d} + 1/\sqrt{m_s} + 1/\sqrt{m_b}\bigr)^2}\,.
  \label{eq:Kinverse}
\end{equation}
At the compositeness scale (which we identify with the adjugate
vacuum), $K^{-1} = 2/3$ exactly.
SM renormalisation group running shifts $K^{-1}$ away from $2/3$
at low energies.

Using the multi-loop Yukawa couplings of Antusch, Hinze, and
Saad~\cite{AHS:2024}, we evaluate $K^{-1}$ as a function of the
RG scale~$Q$ in two scenarios (Table~\ref{tab:RG}).

\begin{table}
\centering\small
\caption{Inverse Koide ratio $K^{-1}(d,s,b)$ under RG running.
  ``Best scale'' is the scale at which $K^{-1}$ is closest
  to~$2/3$.  $\Delta = K^{-1} - 2/3$ in percent.}
\label{tab:RG}
\begin{tabular}{lccr}
\toprule
Scenario & Best scale & $K^{-1}$ & $\Delta$ (\%) \\
\midrule
\textbf{SM desert} & $\mathbf{\sim 10\TeV}$ & $\mathbf{0.66667}$
  & $\mathbf{+0.0003}$ \\
MSSM, $M_S = 3\TeV$, $\tan\beta = 5$ & $M_{\rm GUT}$ & $0.66689$ & $+0.033$ \\
MSSM, $M_S = 3\TeV$, $\tan\beta = 10$ & $M_{\rm GUT}$ & $0.66693$ & $+0.040$ \\
MSSM, $M_S = 3\TeV$, $\tan\beta = 30$ & $M_{\rm GUT}$ & $0.66749$ & $+0.124$ \\
MSSM, $M_S = 3\TeV$, $\tan\beta = 50$ & $M_{\rm GUT}$ & $0.66895$ & $+0.342$ \\
MSSM, $M_S = 10\TeV$, $\tan\beta = 5$  & $M_{\rm GUT}$ & $0.66668$ & $+0.003$ \\
MSSM, $M_S = 10\TeV$, $\tan\beta = 10$ & $M_{\rm GUT}$ & $0.66673$ & $+0.009$ \\
MSSM, $M_S = 10\TeV$, $\tan\beta = 30$ & $M_{\rm GUT}$ & $0.66724$ & $+0.086$ \\
MSSM, $M_S = 10\TeV$, $\tan\beta = 50$ & $M_{\rm GUT}$ & $0.66854$ & $+0.281$ \\
\bottomrule
\end{tabular}
\end{table}

\paragraph{SM desert.}
In the SM with no new thresholds between $M_Z$ and the Planck
scale, $K^{-1}(d,s,b)$ decreases monotonically with $Q$ and
crosses~$2/3$ at $Q \approx 10\TeV$ to a precision
of~$0.0003\%$ using AHS multi-loop
data~\cite{AHS:2024,MartinRobertson:2019}.
At $Q = M_Z$, the deviation is $+0.11\%$: still within~$0.2\%$
of the trace identity, but measurably different from exact
$K^{-1} = 2/3$.

\paragraph{MSSM.}
In the MSSM, $K^{-1}$ at $M_{\rm GUT}$ lies $0.23$--$0.65\%$
above~$2/3$ for all values of $\tan\beta$ and $M_{\rm SUSY}$
at two-loop accuracy~\cite{AHS:2024}.
The one-loop MSSM results in Table~\ref{tab:RG} are somewhat
closer to $2/3$ (missing higher-order effects), but the
qualitative conclusion is robust: \emph{the MSSM never
crosses~$2/3$}.
In no MSSM scenario does $K^{-1}$ reach $2/3$ between
$M_{\rm SUSY}$ and $M_{\rm GUT}$.

\paragraph{Interpretation.}
If the trace identity $K^{-1}(d,s,b) = 2/3$ holds exactly at the
compositeness scale, the RG evidence favours
$\SU(3)_F$ compositeness emerging at $Q \sim 10$--$280\TeV$
in a standard model desert, rather than through a conventional
SUSY threshold.
This does \emph{not} exclude SUSY as such---the sBootstrap itself
relies on $\mathcal{N}{=}1$ holomorphy---but disfavours the MSSM
spectrum as the next threshold above the electroweak scale.


%======================================================================
\section{Discussion}
\label{sec:discussion}
%======================================================================

\subsection{What the duality explains}

The $C_3/C_6$ Dirac--Majorana correspondence provides three
structural insights:

(i)~\emph{Origin of mass-type asymmetry.}
The bilinear/trilinear distinction of $s$-confined composites,
which is an algebraic fact, lifts to the electric/magnetic
distinction of M-theory $p$-form gauge fields.
This is not a coincidence but a consequence of the brane
construction.

(ii)~\emph{Mixing angle hierarchy.}
The K\"ahler protection asymmetry (Sec.~\ref{sec:kahler}) provides
a qualitative explanation for the observed pattern: small CKM
angles (meson sector, diagonal $Z_M$) vs.\ large PMNS angles
(baryon sector, democratic $Z_B$).

(iii)~\emph{Seesaw architecture.}
The superpotential $W = BM\tilde{B}/\Lambda^5$ provides the
type-I seesaw skeleton with baryoninos as right-handed neutrinos
at the charged-lepton mass scale, without introducing a GUT-scale
seesaw.

\subsection{What the duality does not explain}

(i)~\emph{The neutrino mass scale} requires K\"ahler factors of
order $10^5$, which is a statement about the strongly-coupled
K\"ahler potential and cannot be computed from holomorphy.

(ii)~\emph{The PMNS mixing angles} are trivial at tree level on
the mesonic branch and require off-diagonal K\"ahler rotations or
radiative corrections.

(iii)~\emph{A quantitative Dirac--Majorana mass relation} does not
exist within the current framework.
The duality is structural, not predictive for individual masses.

\subsection{Pendleton--Ross fixed point: a debunked coincidence}
\label{sec:pendleton_ross}

The Pendleton--Ross~\cite{PendletonRoss:1981} / Hill~\cite{Hill:1981}
infrared quasi-fixed point for the top Yukawa coupling gives,
at one loop in the QCD-only approximation:
\begin{equation}
  R^* = \frac{y_t^2}{g_3^2} = \frac{2}{9}\,.
  \label{eq:PR}
\end{equation}
This derives from the one-loop $\beta$-function
$16\pi^2\,\dd R/\dd t = g_3^2\,R\,(9R - 2)$,
giving $R^* = 2/9$ at the zero.
The coincidence of $R^* = 2/9$ with the Cartan angle
$\delta = 2/9$ of the sBootstrap was noted as potentially
significant.

We have verified computationally that this coincidence is
\emph{spurious}, for three independent reasons:

(i)~Electroweak corrections shift the quasi-fixed point from
$2/9 \approx 0.222$ to $R^* \approx 0.39$ at $M_Z$---a $76\%$
increase.
The shift is:
\begin{equation}
  R^*(Q)
    = \frac{2}{9}
    + \frac{1}{9}\!\left[
      \frac{17}{10}\frac{g_1^2}{g_3^2}
    + \frac{9}{2}\frac{g_2^2}{g_3^2}
    \right]\!.
  \label{eq:PR_EW}
\end{equation}

(ii)~With running gauge coupling ratios, $R^*$ is
scale-dependent: there is no true fixed point in the full SM.
A numerical scan with initial $R \in [0.1, 5.0]$ at
$Q = 10^{16}\GeV$ shows no convergence to a single value at
$M_Z$~\cite{Zimmermann:1985}.

(iii)~The physical SM ratio is $R = y_t^2/g_3^2 = 0.59$ at
$M_Z$, which is $2.66$ times the QCD-only value.
The QCD-only fixed point predicts
$m_t \approx 100\GeV$, off by $73\%$.

\emph{The Pendleton--Ross ratio $2/9$ is exact only in the
one-loop QCD-only limit, which is unphysical.
The agreement with~$\delta = 2/9$ is a numerical accident.}

\subsection{Experimental tests}

The strongest near-term test of the sBootstrap neutrino sector
is the mass ordering.
The geometric seesaw with normal-ordering Dirac mass structure
predicts $m_1 < m_2 < m_3$.
Current data (DESI + CMB) favour normal ordering with
$\sum m_\nu < 0.057$~eV~\cite{PDG:2024}; the
sBootstrap gives $\sum m_\nu \approx 0.060$~eV~\cite{Rivero:2024epjc},
marginally consistent.
JUNO~\cite{JUNO:2022} will definitively determine the mass
ordering by 2027--2028.
\emph{If JUNO establishes inverted ordering, the sBootstrap
seesaw structure is disfavoured.}

The inverse Koide crossing at $Q \approx 10\TeV$ predicts no new
coloured particles below this scale.
This is testable at the HL-LHC and future colliders: observation
of MSSM-like superpartners at $M_S < 10\TeV$ would conflict with
the SM-desert scenario favoured by the RG analysis.

The qualitative prediction of small CKM vs.\ large PMNS mixing
from the K\"ahler asymmetry could be sharpened by a computation
of the baryon K\"ahler metric using holographic QCD
techniques~\cite{SakaiSugimoto:2005}.


%======================================================================
\section{Conclusions}
\label{sec:conclusions}
%======================================================================

We have shown that the M-theory lift of the sBootstrap---a
confining $\SU(3)_F$ family gauge theory at
$\Nf = \Nc = 3$---contains a structural correspondence between
the Dirac and Majorana mass types and the electromagnetic duality
of $p$-form gauge fields in 11 dimensions.
Mesons (M2-branes, $C_3$) produce Dirac masses; baryons (M5-branes,
$C_6$) produce Majorana masses.
The involution $M_i \leftrightarrow \adj(M)_i$ lifts to the
exchange of $A$- and $B$-periods of the genus-2 Seiberg--Witten
curve.

The composite structure explains why CKM angles are small (meson
K\"ahler metric is diagonal) and PMNS angles are large (baryon
K\"ahler metric is democratic).
The superpotential $W = BM\tilde{B}/\Lambda^5$ provides a type-I
seesaw skeleton with baryonino Majorana masses
$M_R^{ii} = m_i$ (exact from $s$-confinement), but the physical
neutrino mass scale requires non-perturbatively large K\"ahler
suppression that cannot be computed from holomorphy alone.

Multi-loop SM running of down-quark Yukawa couplings shows that
$K^{-1}(d,s,b) = 2/3$ at $Q \approx 10\TeV$ in the SM desert,
while no MSSM scenario reaches this value---discriminating between
$\SU(3)_F$ compositeness at the $10$--$280\TeV$ scale and
conventional low-energy SUSY.

The previously noted coincidence between the Pendleton--Ross
QCD-only fixed point $y_t^2/g_3^2 = 2/9$ and the Cartan angle
$\delta = 2/9$ has been shown to be spurious: electroweak
corrections shift the quasi-fixed point to $R^* \approx 0.39$,
and the physical SM ratio is $R = 0.59$.

The framework makes three testable predictions:
(i)~neutrino normal ordering (testable by JUNO, 2027--2028);
(ii)~no new coloured particles below $\sim 10\TeV$
(testable at HL-LHC and future colliders);
(iii)~a qualitative asymmetry between CKM and PMNS mixing
rooted in the composite structure of the family sector.

The principal open problems are the computation of the
non-perturbative K\"ahler potential for the baryonic sector
(which would predict PMNS angles and the neutrino mass scale),
and the determination of the Cartan angle $\delta = 2/9$ from
the instanton structure of the confining theory.


%======================================================================
\begin{acknowledgments}
The author thanks M.~Porter for sustained theoretical discussions
that identified Seiberg duality and the mesino interpretation as
the mechanism underlying the sBootstrap.
During the preparation of this work the author used Claude
(Anthropic) for algebraic verification, numerical checks,
literature cross-referencing, and drafting assistance.
The author reviewed and edited all content and takes full
responsibility for the publication.
\end{acknowledgments}


%======================================================================
\begin{thebibliography}{99}

\bibitem{Rivero:2005b}
  A.~Rivero,
  arXiv:hep-ph/0510068 (2005).

\bibitem{Rivero:2024epjc}
  A.~Rivero,
  Eur.\ Phys.\ J.\ C \textbf{84}, 1058 (2024),
  arXiv:2407.05397.

\bibitem{Kaplan:1991dc}
  D.B.~Kaplan,
  Nucl.\ Phys.\ B \textbf{365}, 259 (1991).

\bibitem{Seiberg:1994bz}
  N.~Seiberg,
  Nucl.\ Phys.\ B \textbf{435}, 129 (1995),
  arXiv:hep-ph/9411149.

\bibitem{HananyWitten:1997}
  A.~Hanany and E.~Witten,
  Nucl.\ Phys.\ B \textbf{492}, 152 (1997),
  arXiv:hep-th/9611230.

\bibitem{Witten:1997M5}
  E.~Witten,
  Nucl.\ Phys.\ B \textbf{500}, 3 (1997),
  arXiv:hep-th/9703166.

\bibitem{Witten:1998baryon}
  E.~Witten,
  JHEP \textbf{07}, 006 (1998),
  arXiv:hep-th/9805112.

\bibitem{CremmerJuliaScherk:1978}
  E.~Cremmer, B.~Julia, and J.~Scherk,
  Phys.\ Lett.\ B \textbf{76}, 409 (1978).

\bibitem{Townsend:1995}
  P.K.~Townsend,
  Phys.\ Lett.\ B \textbf{350}, 184 (1995),
  arXiv:hep-th/9501068.

\bibitem{CsakiSchmaltzSkiba:1997}
  C.~Cs\'aki, M.~Schmaltz, and W.~Skiba,
  Phys.\ Rev.\ Lett.\ \textbf{78}, 799 (1997),
  arXiv:hep-th/9610139.

\bibitem{SeibergWitten:1994}
  N.~Seiberg and E.~Witten,
  Nucl.\ Phys.\ B \textbf{426}, 19 (1994),
  arXiv:hep-th/9407087.

\bibitem{ArgyresPlesserShapere:1996}
  P.C.~Argyres, M.R.~Plesser, and A.D.~Shapere,
  Nucl.\ Phys.\ B \textbf{483}, 172 (1997),
  arXiv:hep-th/9608129.

\bibitem{HoriOoguriOz:1997}
  K.~Hori, H.~Ooguri, and Y.~Oz,
  Adv.\ Theor.\ Math.\ Phys.\ \textbf{1}, 1 (1997),
  arXiv:hep-th/9706082.

\bibitem{MasieroVeneziano:1985}
  A.~Masiero and G.~Veneziano,
  Nucl.\ Phys.\ B \textbf{249}, 593 (1985).

\bibitem{Karasik:2020}
  A.~Karasik,
  JHEP \textbf{02}, 166 (2021),
  arXiv:2012.11154.

\bibitem{AHS:2024}
  S.~Antusch, C.~Hinze, and S.~Saad,
  arXiv:2510.01312 (2025).

\bibitem{MartinRobertson:2019}
  S.P.~Martin and D.G.~Robertson,
  Phys.\ Rev.\ D \textbf{100}, 073004 (2019),
  arXiv:1907.02500.

\bibitem{PendletonRoss:1981}
  B.~Pendleton and G.G.~Ross,
  Phys.\ Lett.\ B \textbf{98}, 291 (1981).

\bibitem{Hill:1981}
  C.T.~Hill,
  Phys.\ Rev.\ D \textbf{24}, 691 (1981).

\bibitem{Zimmermann:1985}
  W.~Zimmermann,
  Commun.\ Math.\ Phys.\ \textbf{97}, 211 (1985).

\bibitem{PDG:2024}
  S.~Navas \textit{et al.} (Particle Data Group),
  Phys.\ Rev.\ D \textbf{110}, 030001 (2024).

\bibitem{JUNO:2022}
  A.~Abusleme \textit{et al.} (JUNO Collaboration),
  J.\ Phys.\ G \textbf{43}, 030401 (2022),
  arXiv:2204.13249.

\bibitem{SakaiSugimoto:2005}
  T.~Sakai and S.~Sugimoto,
  Prog.\ Theor.\ Phys.\ \textbf{113}, 843 (2005),
  arXiv:hep-th/0412141.

\bibitem{Koide:1983qe}
  Y.~Koide,
  Phys.\ Lett.\ B \textbf{120}, 161 (1983).

\bibitem{LutyTerning:2000}
  M.A.~Luty and J.~Terning,
  Phys.\ Rev.\ D \textbf{62}, 075006 (2000),
  arXiv:hep-ph/9812482.

\bibitem{Strominger:1996}
  A.~Strominger,
  Phys.\ Lett.\ B \textbf{383}, 44 (1996),
  arXiv:hep-th/9512059.

\bibitem{BargerBergerOhmann:1993}
  V.~Barger, M.S.~Berger, and P.~Ohmann,
  Phys.\ Rev.\ D \textbf{47}, 1093 (1993).

\end{thebibliography}

\end{document}
